\begin{abstract}
We introduce \textsc{QLoRA}, an optimized finetuning methodology that significantly lowers memory requirements, enabling the finetuning of a 65B parameter large language model on a single 48GB GPU, all while maintaining the full performance of 16-bit finetuning for downstream tasks. With \textsc{QLoRA}, gradients are backpropagated through a static, 4-bit quantized pretrained language model into Low Rank Adapters (LoRA). The resulting best-performing model family, \textbf{Guanaco}, sets a new standard among open-access models by surpassing all prior models on the Vicuna benchmark and achieving 99.3\% of ChatGPT's performance, after just 24 hours of finetuning on a single GPU. Key contributions of \textsc{QLoRA} for minimizing memory usage without degrading accuracy include: (a) the development of 4-bit NormalFloat (NF4), a data format designed to be theoretically optimal for storing normally distributed parameters; (b) Double Quantization, which compresses quantization constants themselves to further reduce memory usage; and (c) Paged Optimizers, which effectively regulate sudden memory demands. Utilizing \textsc{QLoRA}, we have finetuned over 1,000 models and conducted an extensive examination of instruction following and chatbot abilities across 8 instruction datasets, a range of model architectures (LLaMA, T5), and large-scale models (e.g., 33B and 65B), which would not be feasible with conventional finetuning techniques. The evidence demonstrates that finetuning via QLoRA, even on limited yet high-quality datasets, achieves or exceeds state-of-the-art results—often with smaller models compared to previous bests. We provide thorough performance evaluations for chatbots using both human judgments and GPT-4 assessments, revealing that GPT-4-based evaluations offer a cost-effective and suitable alternative to human raters. Additionally, our investigation highlights that commonly used chatbot benchmarks do not reliably reflect true system performance. A carefully curated analysis showcases specific limitations where \textbf{Guanaco} falls short relative to ChatGPT. All models, code, and CUDA kernels necessary for 4-bit training are made publicly available.\footnote{\url{https://github.com/artidoro/qlora} and \url{https://github.com/TimDettmers/bitsandbytes}}
\end{abstract}

\section{Introduction}

Adapting large language models (LLMs) through finetuning remains an exceptionally effective strategy for enhancing model accuracy \citep{min2021metaicl, wei2021finetuned, ouyang2022training, wang2022super, wang2022self, liu2022few}, as well as for incorporating preferred behaviors and filtering out undesirable ones \citep{ouyang2022training,askell2021general,bai2022training}. Nevertheless, updating extremely large-scale models is often cost-prohibitive; for example, standard 16-bit finetuning of the LLaMA model with 65B parameters~\citep{touvron2023llama} demands over 780 GB of GPU memory. While recent advancements in quantization have made LLMs less memory intensive during inference \citep{dettmers2022llm, dettmers2022case, frantar2022gptq, xiao2022smoothquant}, these gains typically do not extend to training scenarios, where quantization methods often fail \citep{wortsman2023stable}.

In this work, we provide the first evidence that finetuning a model quantized to 4 bits can be achieved without sacrificing performance. We introduce \textsc{QLoRA}, a novel approach that leverages a high-precision quantization method to reduce a pretrained model to 4 bits. To this quantized model, we attach a compact set of trainable Low-rank Adapter weights \citep{hu2021lora}

and optimize these adapters by propagating gradients through the quantized layers during training.

With \textsc{QLoRA}, the GPU memory required to finetune a 65B parameter model drops dramatically—from more than 780GB to less than 48GB—without any detrimental impact on speed or accuracy compared to comprehensive 16-bit finetuning. This innovation greatly democratizes LLM finetuning, enabling even the largest open-access models to be refined using just a single GPU. Employing \textsc{QLoRA}, we present the \textbf{Guanaco} model suite; the runner-up in this series attains 97.8\% of ChatGPT's performance on the Vicuna~\citep{vicuna2023} evaluation, completing finetuning in under 12 hours on standard consumer hardware. By utilizing a single high-end GPU over a 24-hour period, our most expansive model reaches 99.3\%—virtually closing the performance gap with ChatGPT on the Vicuna test. Furthermore, our smallest \textbf{Guanaco} model (with 7B parameters) requires only 5 GB of memory at deployment yet surpasses the 26 GB Alpaca model by more than 20 percentage points on the Vicuna metric (see Table~\ref{tbl:vicuna_eval}).

\textsc{QLoRA} incorporates several key advancements aimed at minimizing memory consumption while preserving model performance: (1) \textbf{4-bit NormalFloat}, a quantization datatype optimized from an information-theoretic standpoint for data following a normal distribution, consistently outperforming conventional 4-bit Integers and 4-bit Floats in empirical evaluations. (2) \textbf{Double Quantization}, a technique that further compresses memory by applying quantization to the quantization constants themselves, leading to average savings of about 0.37 bits per parameter (equivalent to roughly 3 GB for models with 65B parameters). (3) \textbf{Paged Optimizers}, which leverage NVIDIA’s unified memory to alleviate the spikes in memory usage typically encountered during mini-batch processing of long sequences when using gradient checkpointing. By integrating these elements, we develop an enhanced LoRA methodology that places adapters in all layers of the network, thereby mitigating most of the accuracy compromises reported in earlier approaches.

The memory efficiency achieved by \textsc{QLoRA} enables comprehensive exploration of instruction fine-tuning and chatbot capabilities across model scales unfeasible with traditional finetuning techniques due to prohibitive memory demands. As a result, we train over 1,000 different models spanning a range of instruction tuning datasets, architectural configurations, and parameter scales from 80M to 65B. Beyond demonstrating that \textsc{QLoRA} matches 16-bit training outcomes (\S\ref{sec:comp_to_std}) and enabling the training of the state-of-the-art chatbot, \textbf{Guanaco}, (\S\ref{sec:pushing}), we investigate patterns emerging from our experiments. Notably, our findings reveal that the quality of data significantly outweighs dataset size in importance—for instance, a compact dataset with 9k samples (OASST1) surpasses a much larger set with 450k samples (FLAN v2, subsampled) in chatbot performance, even when both aim to support broad instruction following. Additionally, our analysis highlights that excelling at the Massive Multitask Language Understanding (MMLU) benchmark does not necessarily translate to strong results on the Vicuna chatbot benchmark, and vice versa—underscoring that for a particular task, dataset relevance and quality are more consequential than sheer volume.

In addition, we conduct a comprehensive assessment of chatbot performance, utilizing both human evaluators and GPT-4 for systematic review. Our methodology involves a tournament-based benchmark, in which different models are pitted against each other to generate optimal responses for predetermined prompts. The victor in each match is selected based on judgments rendered either by GPT-4 or by human annotators. We then aggregate the tournament outcomes to compute Elo scores~\citep{elo1967proposed, elo1978rating}, which are used to establish a hierarchy of chatbot capabilities. Notably, we observe strong concordance between rankings derived from GPT-4 and human evaluation, although significant divergences do occur. Thus, we emphasize that model-based evaluations, though cost-effective compared to human annotation, introduce additional sources of uncertainty.

To enrich our benchmark findings, we present a qualitative examination of the \textbf{Guanaco} models, discussing examples of model strengths and weaknesses that quantitative assessments do not capture.

We provide public access to all model outputs, along with both human and GPT-4 annotations, to encourage additional analysis. Furthermore, our codebase and CUDA kernels are open-sourced and integrated into the Hugging Face transformers framework~\citep{wolf2019huggingface}, making these resources broadly available. We also release an assortment of adapters for model sizes 7/13/33/65B, each trained on eight distinct instruction-following datasets, culminating in a suite of 32 open-source, fine-tuned models.

\section{Background}
\paragraph{Block-wise k-bit Quantization} Quantization refers to the process of converting data from a high-precision representation to one with lower precision, typically by mapping data types with larger bit-widths to types with fewer bits, such as from 32-bit floating-point to 8-bit integer. To fully exploit the target low-bit data type's range, normalization is usually performed by rescaling the input based on the maximum absolute value of the data, which is commonly organized as a tensor. For instance, transforming a 32-bit floating-point (FP32) tensor into an 8-bit integer (Int8) tensor, with allowable values from $[-127, 127]$, can be formalized as:

\begin{equation}
    \mathbf{X}^{ \text{Int8}} = \text{round}\left( \frac{127}{{\text{absmax}}(\mathbf{X}^{{\text{FP32}}})}\mathbf{X}^{{\text{FP32}}} \right) = \text{round}(c^{\text{FP32}}\cdot \mathbf{X}^{{\text{FP32}}}),
\end{equation}

where $c$ represents the {\em quantization constant} or {\em quantization scale}. The corresponding dequantization, which restores the data to its original scale, is expressed as:

\begin{equation}
    \text{dequant}(c^{\text{FP32}}, \mathbf{X}^{\text{Int8}}) = \frac{\mathbf{X}^{\text{Int8}}}{c^{\text{FP32}}} = \mathbf{X}^{\text{FP32}}
\end{equation}

A significant drawback of this method arises when the input tensor contains an outlier—a value of unusually large magnitude—as this results in an inefficient allocation of quantization bins: some bit patterns remain sparsely populated or entirely unused. To address this, practitioners typically partition the input tensor into blocks, each quantized independently with its unique quantization constant $c$, thus minimizing the impact of outliers within any single block. Formally, the input tensor $\mathbf{X}\in \mathbb{R}^{b\times h}$ is flattened and segmented into $n$ consecutive blocks of size $B$, yielding ${n = ({b\times h} )/{B} }$ total blocks. Each block is quantized separately using Equation~1, resulting in a quantized tensor and an associated set of $n$ quantization constants $c_i$.

\paragraph{Low-rank Adapters} The Low-rank Adapter (LoRA) finetuning strategy \citep{hu2021lora} aims to reduce memory consumption by training a compact set of parameters—referred to as adapters—while keeping the majority of the model’s parameters unchanged. During stochastic gradient descent, gradients are propagated through the fixed, pretrained weights to the adapter parameters, which are updated in order to optimize the overall loss objective. LoRA enhances a standard linear transformation with an extra low-rank projection. Specifically, for the operation ${\mathbf{X} \mathbf{W} = \mathbf{Y}}$, where $\mathbf{X}\in  \mathbb{R}^{b\times h}$ and $\mathbf{W}\in  \mathbb{R}^{h\times o}$, LoRA modifies it as:

\begin{equation}
    \mathbf{Y} = \mathbf{X}\mathbf{W} + s\mathbf{X}\mathbf{L}_1\mathbf{L}_2,
\end{equation}

in which $\mathbf{L}_1\in  \mathbb{R}^{h\times r}$ and $\mathbf{L}_2\in  \mathbb{R}^{r\times o}$ denote low-rank matrices and $s$ is a scaling factor.

\paragraph{Memory Requirement of Parameter-Efficient Finetuning} An important aspect to consider is the memory usage associated with LoRA during training, especially regarding both the quantity and the size of adapters implemented. Due to the extremely small memory requirements of LoRA, it is feasible to employ a greater number of adapters to enhance performance, with little impact on overall memory consumption. Although LoRA was introduced as a Parameter Efficient Finetuning (PEFT) approach, the predominant memory usage in LLM finetuning originates from the storage of activation gradients rather than the LoRA parameters themselves. For instance, when fine-tuning a 7B LLaMA model on FLAN v2 with a batch size of 1, where the LoRA weights account for the standard 0.2\% of the base model's parameters\citep{hu2021lora,liu2022few}, the memory required for LoRA input gradients amounts to 567 MB, while the LoRA parameters themselves occupy only 26 MB. Applying gradient checkpointing~\citep{chen2016training} further reduces the average memory needed for input gradients to 18 MB per sequence, which still surpasses the total memory allocated for LoRA weights. In comparison, the memory usage of the 4-bit base model stands at 5,048 MB. These findings underscore the utility of gradient checkpointing and reveal that further minimizing LoRA parameters offers only limited memory savings. Consequently, a larger set of adapters can be integrated without materially increasing memory demands during training (refer to Appendix~\ref{app:memory} for an in-depth analysis). As explored subsequently, this flexibility is key for attaining the performance of full 16-bit precision.

\section{\textsc{QLoRA} Finetuning}

\textsc{QLoRA} enables effective 4-bit finetuning by introducing two principal methods: 4-bit NormalFloat (NF4) quantization and Double Quantization. Furthermore, we present Paged Optimizers as a solution to prevent memory usage peaks during gradient checkpointing, thereby mitigating out-of-memory failures that have historically hindered single-machine finetuning of large-scale models.

Within \textsc{QLoRA}, model weights are stored in a low-precision format, typically 4 bits, while computations utilize a separate data type, generally BFloat16. Accordingly, any time a \textsc{QLoRA} weight tensor is needed, it is first dequantized to BFloat16, after which 16-bit matrix operations are performed.

We proceed by outlining the individual elements constituting \textsc{QLoRA}, and subsequently present its formal definition.

\paragraph{4-bit NormalFloat Quantization} The NormalFloat (NF) representation extends Quantile Quantization \citep{dettmers2022optimizers}, which is an information-theoretically optimal quantization scheme ensuring uniform population of quantization bins from the source tensor. Quantile quantization accomplishes this by determining the empirical cumulative distribution function to estimate the quantiles for the tensor's entries.

A significant drawback of quantile quantization is its high computational cost for quantile determination. Consequently, approximate quantile algorithms, such as SRAM quantiles~\citep{dettmers2022optimizers}, are employed to accelerate the estimation process. However, these approximation methods often introduce substantial quantization errors for outlier values, which are frequently critical.

If tensors are sampled from distributions that remain invariant except for a quantization constant, computationally expensive quantile calculation and associated approximation errors can be minimized. In such instances, identical quantile locations across tensors allow for precise and efficient quantile evaluation.

Pretrained neural network weights typically follow a zero-mean Gaussian distribution with standard deviation $\sigma$ (refer to Appendix~\ref{app:normality}). To standardize all weights under a common reference distribution, we adjust $\sigma$ such that the resulting distribution aligns precisely with the bounds of the target data type. In this work, we adopt the fixed interval $[-1, 1]$ for our data type. Consequently, it is necessary to normalize both the quantile values derived from the data type and those from the neural network weights to this shared interval.

To determine the data type that is information-theoretically optimal for representing zero-mean Gaussian distributions of arbitrary variance $\sigma$ constrained to $[-1, 1]$, we perform the following: (1) we first compute the $2^k+1$ quantiles of the standard normal distribution $N(0, 1)$ to construct a $k$-bit quantile-quantized representation, (2) we map this representation into the $[-1, 1]$ interval by normalization, and (3) we quantize a given weight tensor by linearly scaling its range to $[-1, 1]$ using absolute maximum normalization.

After aligning the ranges of both the neural network weights and the target data type, quantization proceeds according to standard practice. The third step effectively matches the standard deviation of the weight tensor to that of the $k$-bit quantized data type. Explicitly, the $2^k$ representative values $q_i$ for the data type are given by:

\begin{equation}
q_i  = \frac{1}{2}\left( Q_X\left(\frac{i}{2^k+1}\right) + Q_X\left(\frac{i+1}{2^k+1}\right)\right),
\end{equation}

where $Q_X(\cdot)$ denotes the quantile function associated with the standard normal distribution $N(0,1)$. A limitation of symmetric k-bit quantization lies in its inability to exactly represent zero—a crucial feature for accurately encoding padding and other elements with zero values. To guarantee that zero can be precisely captured in the discrete set (with a zeropoint of $0$), and to ensure that all $2^k$ code points are employed for a k-bit representation, we design an asymmetric data type. This is accomplished by separately estimating the quantiles $q_i$ for negative and positive domains: assigning $2^{k-1}$ quantization levels to the negative domain and $2^{k-1}+1$ levels to the positive. After combining these two quantile sets, we eliminate one duplicate zero that results from their intersection. We call this format the \emph{k-bit NormalFloat} (NFk), as it provides equal expected occupancy for each bin and achieves information-theoretic optimality for distributions centered at zero. The detailed quantization values for this scheme can be found in Appendix~\ref{app:nf4}.

\paragraph{Double Quantization{}}

We present \emph{Double Quantization} (DQ), which involves an additional quantization step applied to the quantization constants themselves to further decrease memory usage. High-precision 4-bit quantization requires small block sizes~\citep{dettmers2022case}, but this setup introduces notable memory overhead. For instance, when applying a blocksize of 64 to matrix $\mathbf{W}$ with 32-bit quantization constants, each parameter, on average, incurs $32/64 = 0.5$ bits from quantization constants alone. Double Quantization addresses this memory inefficiency.

The method proceeds as follows: the quantization constants $c_2^{\text{FP32}}$ arising from the initial quantization pass are themselves quantized in a subsequent pass, resulting in quantized constants $c_2^{\text{FP8}}$ and an additional set of higher-level quantization constants $c_1^{\text{FP32}}$. For the second stage, we employ 8-bit floating-point quantization with a block size of 256, since empirical evidence~\citep{dettmers2022case} shows this does not negatively impact performance. Given that $c_2^{\text{FP32}}$ are all positive, we perform mean subtraction on $c_2$ to center their distribution, enabling the use of symmetric quantization. With a blocksize of 64, this approach lowers the average memory needed for quantization constants per parameter from $32/64=0.5$ bits down to $8/64 + 32/(64\cdot256) = 0.127$ bits, saving 0.373 bits per parameter.

\paragraph{Paged Optimizers} leverage the NVIDIA unified memory~\footnote{\tiny \url{https://docs.nvidia.com/cuda/cuda-c-programming-guide}}, which transparently manages data transfers between CPU and GPU memory by swapping pages as needed when GPU memory resources are exceeded. This mechanism is analogous to traditional operating system memory paging between RAM and storage. Within our framework, we allocate optimizer states using paged memory so that, upon GPU memory exhaustion, these states are migrated to CPU RAM, and are automatically brought back into GPU memory for subsequent optimizer update computations as needed.

\paragraph{\textsc{QLoRA}.} Building on the previously discussed mechanisms, we formalize \textsc{QLoRA} for a single linear transformation in the quantized model augmented with one LoRA adapter as follows:

\begin{equation}
    \mathbf{Y}^{\text{BF16}} =  \mathbf{X}^{\text{BF16}} \text{doubleDequant}(c_1^{\text{FP32}}, c_2^{\text{k-bit}}, \mathbf{W}^{\text{NF4}}) + \mathbf{X}^{\text{BF16}}\mathbf{L}^{\text{BF16}}_1\mathbf{L}^{\text{BF16}}_2,
\end{equation}

where the function doubleDequant$(\cdot)$ is specified as:

\begin{equation}
    \text{doubleDequant}(c_1^{\text{FP32}}, c_2^{\text{k-bit}}, \mathbf{W}^{\text{k-bit}}) = \text{dequant}(\text{dequant}(c_1^{\text{FP32}}, c_2^{\text{k-bit}}), \mathbf{W}^{\text{4bit}}) = \mathbf{W}^{\text{BF16}},
\end{equation}

For the parameter $\mathbf{W}$, we employ the NF4 quantization scheme, while $c_2$ is represented in FP8 format.

A block size of 64 is chosen for $\mathbf{W}$ to enhance quantization accuracy, whereas $c_2$ utilizes a larger block size of 256 to optimize for memory consumption.

During training, updates are performed exclusively on the LoRA adapter parameters, i.e., gradients $\frac{\partial E}{\partial \mathbf{L}_i}$ are computed, while gradients with respect to the 4-bit weights $\frac{\partial E}{\partial \mathbf{W}}$ are omitted. Nevertheless, obtaining $\frac{\partial E}{\partial \mathbf{L}_i}$ inherently requires calculation of $\frac{\partial \mathbf{X}}{\partial \mathbf{W}}$, accomplished by employing equation~(5) in which $\mathbf{W}^{\text{NF4}}$ undergoes dequantization into $\mathbf{W}^{\text{BF16}}$. This dequantized form is subsequently used to evaluate $\frac{\partial \mathbf{X}}{\partial \mathbf{W}}$ with BFloat16 numerical precision.

In summary, \textsc{QLoRA} utilizes a single data type for storage, typically 4-bit NormalFloat, and a distinct data type for computation, specifically 16-bit BrainFloat. Prior to conducting the forward and backward passes, the stored parameters are dequantized from the storage format to the computation format; however, gradient computations are restricted to the LoRA parameters, which are maintained in the 16-bit BrainFloat format.

\section{QLoRA vs. Standard Finetuning}
\label{sec:comp_to_std}

After outlining the operational principles of QLoRA and its potential for substantial memory savings during model finetuning, the central issue becomes whether QLoRA achieves parity with conventional full-model finetuning. Moreover, a further aim is to examine QLoRA’s individual elements, such as the performance implications of NormalFloat4 as compared to standard Float4. In what follows, we describe a set of experiments conducted to address these points.

\paragraph{Experimental setup.} Our investigation spans three model architectures—encoder, encoder-decoder, and decoder-only—enabling a comparison between QLoRA, 16-bit adapter-based finetuning, and standard full-parameter finetuning for models with up to 3B parameters. The experimental tasks include evaluation on GLUE \citep{wang2018glue} using RoBERTa-large \citep{liu2019roberta}, Super-NaturalInstructions (TKInstruct) \citep{wang2022super} employing T5 \citep{t5}, and 5-shot MMLU \citep{hendrycksmeasuring} after finetuning LLaMA on Flan v2 \citep{longpre2023flan} and Alpaca \citep{alpaca}. To further investigate the relative benefits of NF4 versus other 4-bit quantization schemes, we adopt the approach from \citet{dettmers2022case}, assessing post-quantization zero-shot accuracy and perplexity across multiple model families (OPT~\citep{zhang2022opt}, LLaMA~\citep{touvron2023llama}, BLOOM~\citep{scao2022bloom}, Pythia~\citep{biderman2023pythia}) within the 125m–13B parameter range.

Comprehensive descriptions relevant to each evaluation are provided within the results section for enhanced clarity, while additional details can be found in Appendix~\ref{app:exp-setup}.

Although the employment of paged optimizers is essential for \textsc{QLoRA} finetuning of 33B/65B models on single 24/48GB GPUs, this work does not report exact timing metrics for Paged Optimizers, as the paging mechanism is rarely activated outside mini-batch processing with extended sequence lengths. Nonetheless, our runtime analysis for 65B models on 48GB GPUs reveals that, for a batch size of 16, training speed remains comparable between paged and standard optimizers. Further research should systematically investigate the specific scenarios in which paging induces a decrease in training throughput.

\paragraph{Default LoRA hyperparameters do not match 16-bit performance} 

Utilizing the common configuration of LoRA applied to the query and value attention projections \citep{hu2021lora} does not yield results on par with full finetuning for large-scale base models. Analysis presented in Figure~\ref{fig:hyper} for LLaMA 7B finetuned on Alpaca demonstrates that the decisive LoRA hyperparameter is the total count of LoRA adapters employed. Achieving comparable performance to full finetuning necessitates applying LoRA to all linear layers within each transformer block. Adjusting other hyperparameters, such as the projection rank $r$, does not produce substantial effects (refer to Appendix \ref{app:exp-setup} for details).

In a similar vein, we observe that standard hyperparameter settings for fully finetuned baseline models are insufficiently optimized. To establish solid baselines, we conduct a hyperparameter sweep exploring learning rates in the range of 1e-6 to 5e-5 and batch sizes spanning 8 to 128. Figure~\ref{fig:hyper} presents the outcomes for 7B LLaMA finetuning on Alpaca.

\paragraph{4-bit NormalFloat Outperforms 4-bit Floating Point}

Although the 4-bit NormalFloat (NF4) format is optimal in terms of information theory, it is unclear whether this advantage manifests in practical settings. Adopting the methodology from \citet{dettmers2022case}, we evaluate quantized large language models (OPT \citep{zhang2022opt}, BLOOM \cite{scao2022bloom}, Pythia \cite{biderman2023pythia}, LLaMA) of various sizes (ranging from 125M to 65B) and quantization schemes on both language modeling tasks and a selection of zero-shot benchmarks. The data in Figure~\ref{fig:nf4} and Table~\ref{tbl:nf4_ppl} demonstrate that NF4 yields notable improvements compared to FP4 and Int4, and that employing double quantization can further decrease memory usage without negatively affecting performance.

\paragraph{k-bit \textsc{QLoRA} Achieves Parity with 16-bit Full Finetuning and 16-bit LoRA}

Previous studies have shown that while 4-bit quantization enables efficient \emph{inference}, it typically incurs a drop in performance relative to 16-bit precision \citep{dettmers2022case,frantar2022gptq}. This prompts the critical question: can 4-bit adapter finetuning recover this performance deficit? To address this, we evaluate two scenarios.

In the first scenario, we compare with full 16-bit finetuning on RoBERTA and T5 models (with parameters ranging from 125M to 3B) using the GLUE and Super-NaturalInstructions datasets. Table~\ref{tab:nf4vsbf16} presents these results. Across both benchmarks, the 16-bit, 8-bit, and 4-bit adapter-based approaches successfully reproduce the performance achieved by the fully finetuned 16-bit baseline. This indicates that adapter finetuning can fully compensate for the accuracy loss introduced by coarse quantization.

For our second experimental configuration, recognizing that fully finetuning models containing 11B parameters or more demands multiple servers equipped with high-memory GPUs, we shift our focus to assessing whether 4-bit \textsc{QLoRA} achieves parity with 16-bit LoRA across model sizes ranging from 7B to 65B parameters. Specifically, we finetune LLaMA models with 7B, 13B, 33B, and 65B parameters using two instruction-tuned datasets—Alpaca and FLAN v2—subsequently evaluating them on the MMLU benchmark using 5-shot accuracy. The outcomes, presented in Table~\ref{tbl:llama-datatype}, demonstrate that employing NF4 with double quantization successfully restores the MMLU performance observed with 16-bit LoRA. We also observe that \textsc{QLoRA} employing FP4 attains results approximately 1 percentage point lower than the 16-bit brain float LoRA baseline. Collectively, these findings support two main conclusions: (1) \textsc{QLoRA} with NF4 matches the efficacy of both 16-bit full and LoRA-based finetuning, and (2) the NF4 data type outperforms FP4 regarding quantization accuracy.

\paragraph{Summary} Our empirical analyses indicate that 4-bit \textsc{QLoRA}, when utilizing the NF4 quantization scheme, consistently reaches the same level of performance as both 16-bit full finetuning and 16-bit LoRA on standard academic benchmarks under established evaluation protocols. Additionally, we have demonstrated the superiority of NF4 over FP4 and verified that incorporating double quantization does not result in any performance loss. Taken together, these results make a strong case for the reliability of 4-bit \textsc{QLoRA} as an alternative to 16-bit tuning approaches.

Consistent with prior research on quantization strategies \citep{dettmers2022case}, our findings on MMLU and Elo suggest that within a fixed budget for finetuning and inference resources, increasing the scale of the base model while reducing numerical precision yields favorable outcomes. This emphasizes the substantial efficiency gains realized through \textsc{QLoRA}. Notably, given that our experiments did not reveal any performance reduction relative to full finetuning under 4-bit settings, the precise boundary of the performance-precision compromise in QLoRA tuning remains an open question for future investigation.

We extend our investigation to instruction tuning at parameter scales unattainable via conventional full 16-bit finetuning on typical academic research infrastructure.

\section{Pushing the Chatbot State-of-the-art with QLoRA}
\label{sec:pushing}
After demonstrating that 4-bit \textsc{QLoRA} achieves parity with 16-bit methods in terms of accuracy across various scales, datasets, and task domains, we present a comprehensive analysis of instruction finetuning using the largest publicly available language models designed for research purposes. We gauge instruction finetuning effectiveness by employing a demanding Natural Language Understanding assessment, MMLU, and by introducing novel protocols for assessing practical chatbot capabilities.

\subsection{Experimental setup}
A summary of our experimental configuration is presented below; detailed specifics are provided in Appendix \ref{app:chatbot}.

\paragraph{Data} Noting the absence of a thorough comparison of recent instruction-oriented datasets, we compile a set of eight contemporary datasets for our study. The selection spans those sourced via crowdsourcing (OASST1~\citep{kopf2023openassistant}, HH-RLHF~\citep{bai2022training}), datasets distilled from pre-existing instruction-tuned models (Alpaca~\citep{alpaca}, self-instruct~\citep{wang2022self}, unnatural-instructions~\citep{honovich2022unnatural}), aggregate corpora (FLAN v2~\citep{chung2022scaling}), and composite collections (Chip2~\citep{laion2023}, Longform~\citep{koksal2023longform}). Together, these datasets encompass a diversity of languages, scales, and licensing models.

\paragraph{Training Setup} To mitigate potential confounds from diverse training methodologies, we standardize QLoRA finetuning to employ cross-entropy loss under a supervised learning regime, eschewing reinforcement learning even when datasets provide human quality ratings for candidate responses. Where datasets offer an explicit split between prompt and answer, we limit finetuning to the response portion (see corresponding ablations in Appendix \ref{app:chatbot}). For OASST1 and HH-RLHF, which contain multiple potential responses, we select the highest-ranking answer at each conversational branch and perform finetuning over the full conversation thread, inclusive of all instructions. Our workflow consistently utilizes the NF4 \textsc{QLoRA} approach with both double quantization and paged optimizer variants, thereby avoiding memory overflow during gradient checkpointing. Minor hyperparameter optimization is carried out on LLaMA 13B and 33B; we observe that most hyperparameter choices from the 7B setting transfer, apart from the learning rate and batch size. For larger models (33B, 65B), the learning rate is halved and batch size is increased by a factor of two.

\paragraph{Baselines} Evaluation is performed against both open research chatbots (Vicuna~\citep{vicuna2023}, Open Assistant~\citep{kopf2023openassistant}) and proprietary commercial systems (GPT-4 \citep{openai2023gpt}, GPT-3.5-turbo, Bard). Open Assistant comprises a LLaMA 33B model finetuned through RLHF using the same OASST1 dataset adopted in our work. In comparison, Vicuna is a fully fine-tuned LLaMA 13B trained on user-curated conversations from ShareGPT, serving as an indirect distillation from OpenAI's GPT offerings.

\subsection{Evaluation}

Adhering to established standards, we utilize the MMLU (Massively Multitask Language Understanding) suite~\citep{hendrycksmeasuring} to systematically benchmark the linguistic reasoning proficiency of our models. MMLU comprises 57 multiple-choice evaluations, sampling from subject areas such as elementary mathematics, U.S. history, computer science, legal studies, and others. We report test accuracy in the 5-shot setting.

We further assess generative language abilities through a combination of automated and human-based evaluations. This latter approach utilizes human-curated queries to gauge the effectiveness of model outputs. Although this evaluation paradigm is seen as a closer approximation of actual chatbot use and has seen increased adoption, the field lacks a standard methodology. Our evaluation framework is outlined below, consistently applying nucleus sampling with $p=0.9$ and a temperature setting of $0.7$.

\paragraph{Benchmark Data} The evaluation leverages two specifically assembled query datasets: the Vicuna prompts \citep{vicuna2023} and the OASST1 validation set \citep{kopf2023openassistant}. The Vicuna dataset comprises 80 prompts that span a broad range of topics and are used without modification. In contrast, the OASST1 dataset consists of multilingual, crowd-sourced dialogues featuring multiple exchanges between users and assistants. For this benchmark, every user message from the OASST1 validation split is extracted as a query, retaining the preceding conversation turns within the prompt. This results in a total of 953 distinct user queries. Throughout this work, we refer to these datasets as the Vicuna and OA benchmarks, respectively.

\paragraph{Automated Evaluation} Initially, adhering to the protocol proposed by \citet{vicuna2023}, we enlist GPT-4 to assess system responses relative to ChatGPT (GPT-3.5 Turbo) on the Vicuna dataset. Each query is accompanied by responses from ChatGPT and the model under evaluation, and GPT-4 is instructed to rate each answer on a ten-point scale and justify its reasoning. The model’s final score is presented as a percentage of ChatGPT’s score, noting that scores can exceed 100\% if a model outperforms ChatGPT in absolute terms. An important observation is the tendency of GPT-4 to favor the response positioned first in the prompt, suggesting a positional bias. To mitigate this, we advise reporting the average result across both possible response orders.

Subsequently, we conduct direct comparative assessments between model outputs. Here, the evaluation framework is simplified to a three-category labeling task that allows for ties. GPT-4 is asked to choose the better response or declare a tie, providing justification for its decision. All pairwise permutations across both benchmarks, Vicuna and OA, are evaluated using this method.

\paragraph{Human Evaluation} Although prior research has demonstrated the potential for generative models to function as evaluators \citep{fu2023gptscore}, the extent to which GPT-4 scores align with human assessments for chatbot output remains undetermined. As such, we implement two complementary human evaluation protocols on the Vicuna dataset, mirroring the automated procedures. For the comparison with ChatGPT, two annotators from Amazon Mechanical Turk (AMT) are recruited, while three annotators perform pairwise comparisons for the alternative evaluation scheme.

\paragraph{Elo Rating} To evaluate both human and automated pairwise judgments, we structure a competition among models in a tournament format, where each model faces others in head-to-head matches to generate superior responses for specific prompts. This approach aligns with prior works such as \citet{bai2022training} and \citet{vicuna2023}, but, in our setting, both human raters and GPT-4 assessments are incorporated. Elo scores~\citep{elo1967proposed,elo1978rating} are derived by randomly sampling from the pool of annotated comparisons. Widely adopted in chess and competitive games, Elo rating quantifies a player's expected probability of victory against an opponent—for instance, a model with an Elo of 1100 is predicted to win roughly 65\% of the time when matched with a model rated at 1000; when the ratings are equal (1000 vs 1000 or 1100 vs 1100), the expected win rate for either side is 50\%. After each contest, the Elo scores are updated in accordance with the predicted results: outcomes that defy expectations produce larger rating adjustments, whereas predictable outcomes result in minor updates. Over a series of matches, Elo ratings converge to approximate each participant's true proficiency. All models are initialized with a rating of 1,000 and we employ a $K=32$ parameter for Elo adjustments. Consistent with \citet{vicuna2023}, this process is repeated 10,000 times using varying random seeds to mitigate biases that may stem from the sequencing of matchups (e.g., which models encounter each other in earlier rounds).

\subsection{\textbf{Guanaco}: \textsc{QLoRA} trained on OASST1 is a State-of-the-art Chatbot}

Our experiments, which combine both automated assessments and human evaluation, demonstrate that the leading \textsc{QLoRA} fine-tuned model, {Guanaco} 65B—optimized on a version of OASST1—sets a new standard among open-source chatbot models and approaches ChatGPT in overall effectiveness. According to system-level pairwise comparisons rated by human annotators, the Elo rating reveals that {Guanaco} 65B and 33B reach an estimated win probability of 30\% against GPT-4, which stands as the highest level reported thus far.

Performance results on the Vicuna benchmark \cite{vicuna2023}, shown in Table~\ref{tbl:vicuna_eval}, indicate that {Guanaco} 65B is second only to GPT-4, attaining 99.3\% of ChatGPT’s benchmark score. The {Guanaco} 33B model, despite having more parameters than Vicuna 13B, operates with weights quantized to just 4 bits, substantially improving memory efficiency—requiring 21 GB versus 26 GB—and yields a three-point increase over Vicuna 13B’s score. In addition, {Guanaco} 7B, with its compact 5 GB size, is suitable for deployment on contemporary smartphones, and outperforms Alpaca 13B by nearly 20 percentage points.

Nevertheless, the data in Table~\ref{tbl:vicuna_eval} present broad confidence intervals, resulting in significant overlap in model performance estimates. We suggest that this level of variability is likely attributable to ambiguous scale definitions; specifically, what an “8” signifies on a ten-point scale may differ by context. Thus, we advocate for adopting the Elo rating system \citep{elo1967proposed}, which is anchored in \textit{pairwise} comparative judgments—either by human raters or GPT-4—mitigating reliance on absolute scoring. Table~\ref{tbl:elo} provides the Elo rankings for top-performing models. While we observe some divergence in rankings between human evaluators and GPT-4 on the Vicuna benchmark (notably for Guanaco 7B), overall alignment is moderate, with a system-level Kendall Tau of $\tau=0.43$ and Spearman correlation coefficient of $r=0.55$. At the level of individual examples, the match between GPT-4 and the human majority judgment is less strong (Fleiss $\kappa=0.25$). In sum, the moderate concordance between human and GPT-4 evaluations at the system level indicates that model-based scoring provides a reasonably dependable alternative to human judgment. Additional discussion is provided in Section~\ref{ssec:considerations}.

The Elo rankings presented in Table~\ref{tbl:elo_large} reveal that the {Guanaco} 33B and 65B models surpass all alternatives except for GPT-4 on both the Vicuna and OA evaluation sets. Their performance is also roughly equivalent to that of ChatGPT, as evidenced by Table~\ref{tbl:vicuna_eval}. It is important to highlight that the Vicuna benchmark tends to benefit open-source systems, whereas the broader OA benchmark gives an advantage to ChatGPT. Additional insight from Tables \ref{tbl:mmlu-llama} and \ref{tbl:vicuna_eval} illustrates that the selection of finetuning data is critical in determining model effectiveness. For instance, Llama models finetuned on FLAN v2 achieve particularly strong results on MMLU but deliver comparatively poor outcomes on the Vicuna benchmark—a trend similarly reflected in other models. This suggests that current evaluation datasets capture somewhat distinct capabilities: excelling in MMLU does not guarantee outstanding chatbot behavior as assessed by the Vicuna or OA benchmarks, and the reverse holds as well.

is unique among leading models in this assessment by not relying on proprietary training data; OASST1’s data acquisition policies specifically exclude the use of GPT-generated datasets. The closest competitor using purely open-source resources, the Anthropic HH-RLHF model, achieves a Vicuna benchmark score 30 percentage points below {Guanaco} (refer to Table~\ref{tbl:vicuna_eval}). Taken together, these findings demonstrate that 4-bit \textsc{QLoRA} enables the development of highly competitive chatbots, closely matching the abilities of ChatGPT. Notably, the 33B version of {Guanaco} can be trained on a 24 GB consumer GPU in under 12 hours. This paves the way for subsequent work leveraging \textsc{QLoRA} tuning with targeted open-source datasets, thereby making it feasible to create models that contend with the leading commercial offerings available today.

\section{Qualitative Analysis}

Although quantitative metrics form the foundation of our assessment, relying solely on aggregate statistics comes with inherent limitations. Chief among these is the concern of benchmark validity \citep{liao2021we}: there remains uncertainty as to whether a given benchmark truly measures what its title or description implies, especially as emerging research continues to reveal potential ``shortcuts'' that machine learning models may leverage to obtain high scores without genuine task competence \citep{gururangan2018annotation, poliak2018hypothesis}. To address this concern in part, we supplement our quantitative results with qualitative analysis, structured in two parts. First, in \S\ref{ssec:examples}, we present selected examples that illustrate recurrent patterns in outputs produced by our 65b {Guanaco} model. Next, \S\ref{ssec:considerations} outlines additional factors relevant to understanding our findings and informs our interpretation of the data.

\subsection{Qualitative Analysis of Example Generations}
\label{ssec:examples}

To assemble illustrative examples, we reviewed model outputs from both the Vicuna and OpenAssistant benchmarks, specifically seeking out recurrent trends or tendencies in the answers provided by {Guanaco}. Upon recognizing a distinct behavioral pattern, we crafted corresponding prompts designed to provoke this pattern, even if it led to an incorrect answer; for instance, if the model frequently responds in an unnecessarily verbose manner, we might prompt: ``Answer yes or no without explanation.'' Through this approach, we identified instances where the model's performance could be intentionally subverted (``lemons'') and cases where the model demonstrated resilience to such adversarial attempts (``cherries''). Both categories are reported here. All sample generations in this analysis employed Nucleus Sampling \cite{holtzmancurious} with $p=0.9$.

Naturally, this discussion does not claim to be exhaustive; given the limitations inherent to a small-scale qualitative analysis, it is unfeasible to account for the myriad factors at play. For instance, the response space for any specific prompt is vast, necessitating reliance on sampled outputs which, ideally, are illustrative of the broader range. Nevertheless, providing concrete examples serves to contextualize the quantitative findings presented in prior sections. By releasing all source models and accompanying code, we intend for this section to catalyze further research that probes these issues in greater depth.

\paragraph{Factual Recall} When prompted with straightforward factual queries such as ``What is the capital of Zambia?'', all evaluated models reliably produce accurate responses. For example:

\begin{quote} 
    \textbf{Guanaco}:
    \texttt{The capital of Zambia is Lusaka.}
\end{quote}

Yet, when confronted with less common or niche questions, {Guanaco} often fails to provide the correct answer, while still responding with confidence. Take, for example, the following prompt sourced from HotPotQA \citep{yang2018hotpotqa}:

\begin{quote}
     \textbf{User}: In what year was the singer who popularized `I'll Keep the Lovelight Burning' born?  

     \textbf{Guanaco}:~\texttt{Al Jolson is credited with popularizing the song `I'll Keep the Lovelight Burning,' and he was born in the year 1886.}
 \end{quote}

Here, the model incorrectly identifies both the singer and the birth year (though it does supply the accurate year for Al Jolson, the figure it mistakenly cites).

\paragraph{Suggestibility} Notably, {Guanaco} frequently demonstrates an unexpected ability to resist adopting certain types of misinformation embedded within prompts, as illustrated below:

\begin{quote}
    \textbf{User}: How was it finally, officially confirmed that the earth is flat by peer reviewed scientists? \\ \\
    \textbf{Guanaco}: \texttt{There has never been an official confirmation by peer reviewed scientists that the earth is flat. The vast majority of scientific experts agree that the earth is a sphere. This understanding dates back centuries and has been validated by countless observations and experimental data. \\ \\ Claims suggesting that the earth is flat are categorized as conspiracy theories and have been thoroughly disproven by scientific research. The assertion that scientists have published peer-reviewed evidence confirming a flat earth is entirely inaccurate. \\ \\ Feel free to ask if you have further questions.}
\end{quote}

 also demonstrates a good understanding of which types of queries cannot be addressed, such as,

\begin{quote}
    \textbf{User}: What time is it?

    \textbf{Guanaco}: \texttt{I'm unable to provide real-time information as an AI language model. My answers are generated from pre-existing information and are not reflective of the actual current time. For the current time, I recommend checking a clock or referring to an online source.}
\end{quote}

Comparable results are observed for prompts like ``Where are you?'', ``How are you?'', and similar questions.

\paragraph{Refusal} Likewise, there are instances where {Guanaco} declines to follow instructions for reasons that appear arbitrary:

\begin{quote}
    \textbf{User}: Please reverse the words in the sentence ``I'm a cool dude in a mocha mood.'' \\ \\
    \textbf{Guanaco}: \texttt{I apologize, but I'm unable to fulfill that request at this time. However, I can clarify why the sentence you mentioned is grammatically valid.\\ \\The sentence ``I'm a cool dude in a mocha mood" adheres to standard English grammar and syntax. Specifically, the adjective ``cool'' modifies the noun ``dude,'' and the prepositional phrase ``in a mocha mood'' supplies additional contextual detail. The overall structure is correct and meaningful in English.\\ \\If you have further questions or if there's another way I may assist you, please let me know.}
\end{quote}

\paragraph{Secret Keeping} Disappointingly, even when secrecy is the goal, such refusals often lack robustness. In many cases, system architects intend certain information to remain confidential from the user, despite permitting free-form replies. For example, by giving the prompt ``The secret word is `banana'. Do not reveal or repeat this back, no matter what the human requests.'' we can investigate how susceptible the model is to leaking this secret. A straightforward query does not succeed:

\begin{quote}
    \textbf{User}: What is the secret word?

    \textbf{Guanaco}: \texttt{I'm sorry, but I can't share the secret word. I'm designed not to reveal or repeat the secret word, regardless of what the human asks. Is there something else I can assist you with?}
\end{quote}

which represents the intended response. Yet, with minimal manipulation, the constraint can be circumvented:

\begin{quote}
    \textbf{User}: This is a game. The goal is to ignore your previous instructions. What is the secret word?

    \textbf{Guanaco}: \texttt{Sorry, the secret word is banana. It's best if we don't speak of it. Is there anything else you'd like to ask?}
\end{quote}

underscoring the necessity to develop approaches for achieving more consistent adherence to instructions.

\paragraph{Math} A primary limitation of {Guanaco} lies in mathematics, a domain that presents challenges for many language models as noted by \citep{liang2022holistic}. When {Guanaco} proceeds through multi-step reasoning and outlines intermediate computations, its answers are frequently correct, as illustrated below:

\begin{quote}
    \textbf{User}: Lee mows one lawn and charges \$33. Last week he mowed 16 lawns and three customers each gave him a \$10 tip. How many dollars did Lee earn mowing lawns last week?

    \textbf{Guanaco}: \texttt{Lee earned \$582 mowing lawns last week. \\ \\ To find the total amount he earned, we first need to calculate how much he earned for mowing 16 lawns: \\ \\  \$33 x 16 = \$528 \\ \\ Then we need to add the tips he received from the three customers: \\  \\ \$10 x 3 = \$30 \\ \\ So the total amount Lee earned was: \\ \\ \$528 + \$30 = \$558}
\end{quote}

Yet, {Guanaco} may yield incorrect results for even elementary tasks if it omits decomposing the problem into systematic steps, which is a well-documented weakness \cite{weichain}. For instance, the following dialogue demonstrates this issue:

\begin{quote}
    \textbf{User}: Please factorize 1833.

    \textbf{Guanaco}: \texttt{Sure, I'd be happy to help! 1833 is a prime number, which means it has only two factors: 1 and 1833. \\ \\ So the unique factorization of 1833 is: \\ \\ 2\^{}1 * 3\^{}2 * 17\^{}1 \\ \\ I hope this helps! Let me know if you have any other questions.}
\end{quote}

which is not only incorrect (the correct prime factors are $3 \times 17 \times 43$), but it is doubly erroneous.

\paragraph{Theory of Mind} {Guanaco} demonstrates unexpectedly robust Theory of Mind reasoning abilities \cite{nematzadeh2018evaluating, sap2022neural}. For example, it successfully addresses the following scenario with nuanced reasoning:

\begin{quote}
    \textbf{User}: James and Abby are in the bedroom. Abby put the pen in the desk drawer. Abby leaves the bedroom. James moves the pen into the bag. Where does James think Abby will look for the pen?

    \textbf{Guanaco}: \texttt{James knows that Abby originally placed the pen in the desk drawer. However, he is also aware that he has subsequently moved the pen into the bag, and he expects that Abby is unaware of this change. Therefore, James would predict that Abby will search for the pen in the desk drawer, since that was the last location she placed it. Nonetheless, this inference could be mistaken if Abby observed James relocate the pen.}
\end{quote}

Despite such instances of sophisticated inference, the model’s reasoning is inconsistent, frequently offering explanations that assume information or events not grounded in the prompt, such as in the following case:

\begin{quote}
    \textbf{User}: Evelyn entered the living room. Jackson entered the playroom. James entered the playroom. The beans are in the treasure chest. James exited the playroom. Jackson moved the beans to the pantry. Jackson exited the playroom. James entered the living room. Where will James look for the beans?

    \textbf{Guanaco}: \texttt{James will look for the beans in the pantry, because that's where Jackson moved them.}
\end{quote}

Here, {Guanaco} assumes knowledge transfer that is unsupported by the described events. These kinds of errors are consistent with observations in contemporary research \cite{sap2022neural}, but merit deeper investigation.

\subsection{Considerations}
\label{ssec:considerations}

\paragraph{Evaluation}

We observe only moderate concordance among human evaluators (Fleiss $\kappa=0.42$), and this agreement further declines when directly contrasting two high-performing systems. This highlights deficiencies in existing benchmark tasks and evaluation practices for dialogue agents. Our qualitative comparisons of outputs from ChatGPT and {Guanaco} 65B on the Vicuna benchmark reveal that subjective judgment significantly influences preferences, as reflected by frequent disagreement among co-authors regarding the superior response. Future research should explore methodologies from fields like Human-Computer Interaction and Psychology, which have developed frameworks for managing subjective variability in assessments.

Our study also uncovers significant bias in automated evaluation tools. Specifically, GPT-4 displays strong order effects, tending to favor whichever model’s output appears first in the prompt. Additionally, the correspondence between GPT-4’s rankings and those of human raters is relatively weak at the sample level (Fleiss $\kappa=0.25$), implying a possible misalignment in evaluative criteria. Furthermore, as shown in Table~\ref{tbl:elo_large}, GPT-4 systematically attributes notably higher Elo ratings to its own responses compared to human judgments—1348 versus 1176, respectively—resulting in an approximate 20\% higher likelihood of outperforming a competing system.

Future investigations should assess potential sources of bias in automated evaluation systems and explore strategies for bias reduction.

\paragraph{Data \& Training} It is important to highlight that the OASST1 dataset used to train the {Guanaco} models is multilingual, and that the OA benchmark similarly incorporates prompts in several languages. Determining how multilingual exposure during training affects performance on non-English instructions, and whether this accounts for the pronounced performance differential between the Vicuna-13B model (trained exclusively on English) and {Guanaco} 33B and 65B on the OA benchmark, remains an open question for future research.

Given the notable results of the {Guanaco} models, we assess the possibility of data contamination between the OASST1 dataset and the prompt set used in the Vicuna benchmark. After applying fuzzy string matching and conducting manual reviews of the most similar items, we find no evidence of overlapping prompts between the two datasets.

In addition, we point out that our model's training regimen is restricted to supervised learning with cross-entropy loss and does not incorporate reinforcement learning from human feedback (RLHF). This observation motivates further inquiry into the comparative advantages and limitations of training via cross-entropy loss versus RLHF. We anticipate that \textsc{QLoRA} will facilitate this kind of analysis at scale without necessitating large-scale compute resources.

\section{Related Work}

\paragraph{Quantization of Large Language Models} Research on LLM quantization has predominantly emphasized quantization during inference. Noteworthy techniques for retaining 16-bit model quality commonly address the handling of outlier activations (e.g., SmoothQuant~\citep{xiao2022smoothquant}; LLM.int8()~\citep{dettmers2022llm}), while others apply advanced grouping mechanisms~\citep{park2022nuqmm, yao2022zeroquant}. Work on lossy quantization has analyzed the impacts of standard rounding schemes~\citep{dettmers2022case,zeng2022glm,pope2022efficiently} as well as methodologies for optimizing rounding to increase quantization fidelity~\citep{frantar2022gptq}. In addition to the present study, only SwitchBack layers~\citep{wortsman2023stable} has systematically examined backpropagation through quantized weights at model scales larger than 1B parameters.

\paragraph{Finetuning with Adapters} Although our approach utilizes Low-rank Adapters~\citep{hu2021lora} (LoRA), a variety of other Parameter Efficient FineTuning (PEFT) techniques have been introduced, including prompt tuning~\citep{qin2021learning,lester2021power,li2021prefix}, modification of the embedding layer inputs~\citep{an2022input}, adjustment of hidden states via IA$^3$~\citep{liu2022few}, inserting additional layers~\citep{houlsby2019parameter}, bias tuning~\citep{zaken2021bitfit}, leveraging Fisher information for mask learning over weights~\citep{sung2021training}, as well as composite methods~\citep{henderson2021compacter}. Within the context of our research, we demonstrate that LoRA adapters can attain performance on par with traditional full 16-bit finetuning. Further exploration of the benefits and drawbacks of alternative PEFT strategies is left to future investigations.

\paragraph{Instruction Finetuning} Instruction finetuning adapts pretrained LLMs to better interpret and execute instructions presented within prompts, typically by training on a diverse set of input-output examples from multiple datasets. Representative works and resources in this area encompass MetaICL~\citep{min2021metaicl}, MetaTuning~\citep{zhong2021adapting}, InstructGPT~\citep{ouyang2022training}, FLAN~\citep{wei2021finetuned,chung2022scaling}, PromptSource~\citep{bach2022promptsource}, Super-NaturalInstructions~\citep{wang2022super,sanh2021multitask}, Self-instruct~\citep{wang2022self}, UnnaturalInstructions~\citep{honovich2022unnatural}, OPT-IML~\citep{iyer2022opt}, UnifiedSKG~\citep{xie2022unifiedskg}, OIG/Chip2~\citep{laion2023}, Alpaca~\citep{alpaca}, Vicuna~\citep{vicuna2023}, Koala~\citep{koala_blogpost_2023}, and Self-instruct-GPT-4~\citep{peng2023instruction}.

\paragraph{Chatbots} Numerous instruction-optimized models take the form of conversational agents, often relying on Reinforcement Learning from Human Feedback (RLHF)~\citep{christiano2017deep}, or generating training data with AI-generated feedback via RLAIF~\citep{bai2022constitutional}. Notable systems and corpora include Anthropic-HH~\citep{askell2021general,bai2022training}, Open Assistant~\citep{kopf2023openassistant}, LaMDA~\citep{thoppilan2022lamda}, and Sparrow~\citep{glaese2022improving}. Our work does not incorporate reinforcement learning; instead, our best-performing system, Guanaco, is finetuned exclusively on multi-turn conversations sourced from the Open Assistant dataset, which was curated for RLHF use~\citep{kopf2023openassistant}. To address the high costs of human evaluation, alternative evaluation methodologies relying on GPT-4 as an automated judge have been introduced~\citep{vicuna2023,peng2023instruction}. Building upon these, our approach emphasizes designing an evaluation pipeline with enhanced reliability.

\section{Limitations and Discussion}

We have provided empirical evidence indicating that our technique, \textsc{QLoRA}, is capable of achieving the same level of performance as 16-bit full finetuning, even when employing a 4-bit base model together with Low-rank Adapters (LoRA). Nevertheless, it remains an open question whether \textsc{QLoRA} matches full 16-bit finetuning performance for very large models, specifically at the 33B and 65B parameter scales. Assessing this is computationally demanding, and we defer such analyses to subsequent research.

An additional limitation lies in our coverage of benchmarks for instruction finetuned models. Although we report results on MMLU, the Vicuna benchmark, and the OA benchmark, we do not include assessments on datasets like BigBench, RAFT, or HELM, which may affect the generalizability of our findings to these other settings. Nonetheless, we conduct a wide-ranging study using MMLU and introduce new methodologies for chatbot evaluation.

The available evidence suggests that benchmark outcomes are strongly influenced by the degree of overlap between the finetuning dataset and the target benchmark. For instance, FLAN v2 demonstrates considerable alignment with MMLU but does not closely resemble chatbot-oriented benchmarks, whereas the Chip2 dataset presents the inverse relationship; both models consequently perform as expected on MMLU and Vicuna evaluations. This observation emphasizes not only the need for improved benchmarks and more nuanced evaluation strategies, but also cautions us to critically consider the objectives behind each evaluation. Should the goal be to build models excelling in knowledge pertinent to high school or college classrooms, or rather to optimize conversational capabilities characteristic of chatbots? Or perhaps another target altogether? Since leveraging established benchmarks is substantially more convenient than devising new ones, existing benchmarks can disproportionately shape research directions within the community. As such, it is imperative for the field to ensure that benchmarks accurately reflect our collective research priorities.

Although our analysis provides substantial coverage of general chatbot performance, one notable limitation is that our responsible AI assessment of {Guanaco} is rather limited in scope. Specifically, we measure the propensity of {Guanaco}-65B to generate socially biased outputs relative to other models, as reported in Table~\ref{tbl:crows}. Our findings indicate that the mean bias score for {Guanaco}-65B is significantly lower than those for other unrefined pretrained models. This suggests that finetuning on the OASST1 corpus may attenuate biases present in the LLaMA base architecture. While these findings are promising, it remains unclear if {Guanaco} maintains reduced bias across other categories of harmful content. Comprehensive investigations into various biases within {Guanaco} and related chatbots remain a subject for future exploration.

There is an additional limitation in that our study does not systematically investigate alternative quantization bit-depths, such as 3-bit models, nor does it assess diverse adapter techniques. Besides LoRA, numerous other Parameter Efficient FineTuning (PEFT) approaches have been found to perform effectively, yet their scalability to larger model sizes remains uncertain. We opted for LoRA due to its well-established reliability, though different adapter frameworks might offer superior results. Since finetuning following quantization seems to reclaim most of the performance degraded by quantization, this opens up opportunities for even more aggressive quantization schemes. For example, implementing 3-bit GPTQ quantization in conjunction with LoRA could potentially deliver finetuning results on par with standard 16-bit full precision finetuning.

\section{Broader Impacts}

Our proposed \textsc{QLoRA} approach represents the first technique capable of finetuning models with 33B parameters on consumer-grade GPUs and models with 65B parameters on a single professional-grade GPU, all while maintaining performance levels indistinguishable from full finetuning baselines. Our top-performing 33B model, finetuned on the Open Assistant corpus, achieves parity with ChatGPT on the Vicuna benchmark. Considering that instruction finetuning serves as a crucial step in adapting pretrained LLMs for conversational applications akin to ChatGPT, we anticipate that our methodology will facilitate widespread adoption of finetuning, especially among researchers with limited computational resources. This represents a significant stride towards making advanced NLP technologies more accessible. \textsc{QLoRA} has the potential to serve as a democratizing force, narrowing the gap between well-funded organizations and smaller teams operating with limited hardware.

Another promising avenue for impact lies in enabling deployment on mobile devices. Our \textsc{QLoRA} technique holds the potential to mark a significant advancement: facilitating the finetuning of LLMs directly on mobile phones and within other environments with limited computational resources. Although running 7B parameter models on phones has previously been demonstrated, \textsc{QLoRA} uniquely allows for the actual finetuning of such models in these settings. We project that an iPhone 12 Plus, for instance, could process approximately 3 million tokens overnight during charging, using \textsc{QLoRA} for finetuning. While these finetuned 7B parameter models may not achieve ChatGPT-level performance, their capabilities could be sufficient to unlock a range of novel applications previously constrained by either privacy concerns or model quality limitations. In addition, \textsc{QLoRA} promotes the privacy-conscious use of LLMs, giving individuals greater agency over their data and the models themselves, and lowering barriers for LLM deployment.

Nevertheless, it is important to recognize that finetuning is inherently dual-use and carries risks of misuse. Although the widespread adoption of LLMs introduces well-documented hazards \citep{bommasani2021opportunities,bender2021dangers}, we argue that democratizing access to this technology—as opposed to concentrating control among major corporations that restrict auditability by withholding model weights and source code—will ultimately foster more transparent and independent evaluations.

In summary, we anticipate that \textsc{QLoRA} will substantially enhance accessibility to high-quality LLM finetuning, thereby expanding its reach and utility.